% Paper 3 — Kaelion OTOC KB
% Empirical Characterization of the Accessibility Flow
% E. F. Perez-Eugenio, 2026
% Format: PRL (revtex4-2, two-column)

\documentclass[aps,prl,twocolumn,showpacs,preprintnumbers,superscriptaddress]{revtex4-2}

\usepackage{amsmath,amssymb}
\usepackage{graphicx}
\usepackage{dcolumn}
\usepackage{bm}
\usepackage{hyperref}

\begin{document}

\title{Empirical Characterization of an Autonomous Accessibility Flow:\\
A Universal $\beta$-Function and Continuous IR Fixed-Point Spectrum\\
from Out-of-Time-Order Correlators}

\author{E.~F.~Perez-Eugenio}
\affiliation{Independent Researcher, Le\'on, Guanajuato, Mexico}
\affiliation{ORCID: 0009-0006-3228-4847}

\date{\today}

% ─────────────────────────────────────────────────────────────────────────────
\begin{abstract}
We report the first systematic empirical characterization of an autonomous
accessibility flow extracted from out-of-time-ordered correlators (OTOCs)
across eleven quantum many-body models spanning four dynamical classes.
The effective $\beta$-function $\beta_{\rm eff}(\Omega) \equiv
d\Omega/d(\ln N)$, where $\Omega_{\rm late}$ is the late-time OTOC
saturation value, collapses onto the universal logistic form
$\beta = -c\,\Omega(1-\Omega)$ across models
(slope $= 1.003$, $\chi^2/\text{dof} = 1.11$, $p = 0.31$),
demonstrating that the flow is \emph{autonomous}---depending only on the
current value of $\Omega$ and not on the system size $N$ at which it is
measured.
The scaling exponent $c$ provides a continuous, observable-independent
diagnostic of the effective infrared (IR) fixed point, spanning
$c \in [0.17, 4.30]$ from integrable-adjacent (Luttinger class, $c \approx
1$) to maximal scramblers (Schwarzian/JT class, $c \approx 3.9$--$4.3$).
A pre-registered prediction---$c \in [1, 4]$ for intermediate-range
interactions---is confirmed with $c = 3.41 \pm 0.48$ for power-law coupling
$J(r) = J_0/r$.
The exponent $c$ is invariant under change of local observable ($1.2\%$
variation between $V = Z_1$ and $V = Z_2$) and UV symmetry class ($4\%$
variation across Altland-Zirnbauer classes in SYK), establishing it as an
infrared property insensitive to ultraviolet structure---the standard
signature of renormalization group universality.
\end{abstract}

\maketitle

% ─────────────────────────────────────────────────────────────────────────────
\section{Introduction}
\label{sec:intro}

The out-of-time-ordered correlator (OTOC)~\cite{Larkin1969,Maldacena2016}
has emerged as a central diagnostic of quantum information scrambling,
probing how initially local operators spread across the degrees of freedom
of a many-body system~\cite{Swingle2018,XuSwingle2024}.
Early-time exponential growth characterizes chaotic dynamics through the
Lyapunov exponent~\cite{Maldacena2016}, while late-time behavior encodes
the system's approach to a fully scrambled state.

The late-time saturation value of the OTOC carries information beyond the
Lyapunov regime.
Huang, Brand\~ao, and Zhang~\cite{HuangFiniteSize2019} established that
in systems with energy conservation, this saturation value scales
polynomially with system size $N$ rather than exponentially.
In Paper~I of this series~\cite{PerezEugenio2026a}, we introduced
$\Omega_{\rm late} \equiv \langle |C(d)| \rangle_{d \geq D_{\rm max}/2}
/ C(0)$ as a fit-free diagnostic that correctly classifies dynamical
regimes---integrable, mixing, chaotic, and disorder-averaged---where
Lyapunov-based measures fail at finite $N$ due to finite-size recurrences.
Paper~II~\cite{PerezEugenio2026b} extended this diagnostic to spatial scrambling
profiles and cross-platform hardware validation, demonstrating that
$\Omega_{\rm late}$ is a hardware-independent observable and justifying
its use as the central quantity of the present analysis.

Two questions motivate the present work.
First, how does $\Omega_{\rm late}(N)$ flow with system size, and is that
flow universal?
Second, can the functional form of the flow classify models according to
their dynamical universality class?

We show that the answer to both questions is yes.
The effective $\beta$-function $\beta_{\rm eff} \equiv d\Omega/d(\ln N)$
is autonomous---it depends only on the current value of $\Omega$, not on
the system size $N$---and collapses onto the logistic form
$\beta = -c\,\Omega(1-\Omega)$ with a single model-dependent exponent $c$.
This is the defining property of a Wilsonian renormalization group (RG)
flow~\cite{Kim2025,Vidal2007,Swingle2012}.
The coefficient $c$ serves as a continuous, observable-independent
spectrometer of the effective IR fixed point, resolving eleven models into
a structured spectrum that reflects their underlying dynamical symmetry.

% ─────────────────────────────────────────────────────────────────────────────
\section{Setup}
\label{sec:setup}

\subsection{Models and Parameters}

We study eleven models spanning four dynamical classes, all using
parameter set B (Conjunto~B): open boundary conditions (OBC),
$J = 1.0$, $h = 0.5$, $b = 0.5$ for the kicked Ising (KI) chaotic
regime, as defined in Paper~I~\cite{PerezEugenio2026a}.
The models are:

\emph{KI chaotic} ($J = 0.6$--$1.0$): Floquet operator
$U_F = e^{-ih\sum X_i} e^{-ib\sum Y_i} e^{-iJ\sum Z_i Z_{i+1}}$
with OBC, varied coupling $J$.

\emph{KI power-law coupling}: $J(r) = J_0/r^\alpha$ ($J_0 = 1.0$,
$\alpha = 1.0$), all-to-all pairs $i < j$.

\emph{KI mixing}: Same Hamiltonian with $J = 0.1$, non-chaotic mixing
regime~\cite{GarciaMata2026,Znidaric2024}.

\emph{SYK-inspired} ($q = 4, 6$): Random all-to-all ZZ couplings
$J_{ij} \sim \mathcal{U}(0.5, 1.5)$, Floquet operator
$U_{\rm SYK} = e^{-iH_{\rm SYK}}$, disorder-averaged over $10$ seeds.

\subsection{Observable and Scaling Analysis}

The OTOC signal uses the forward-backward echo protocol of Paper~I,
\begin{equation}
  C(d) = \langle \phi_r(d) | Z_0 | \phi_r(d) \rangle, \quad
  |\phi_r(d)\rangle = U_F^{-d} Z_r U_F^d |0\rangle^{\otimes N},
\end{equation}
with $r = 1$ (nearest neighbor) unless otherwise noted.
The late-time scrambling residual is
\begin{equation}
  \Omega_{\rm late}(N) = \frac{\langle |C(d)| \rangle_{d \geq D_{\rm max}/2}}{C(0)},
\end{equation}
where $D_{\rm max} = 5N$ ensures full equilibration.

\emph{Methodological note:} For KI chaotic $J = 1.0$ ($c \approx 1$,
slow scrambling), $D_{\rm max} = 3N$ produces a $64\%$ artifact in
$\Omega_{\rm late}$.
$D_{\rm max} = 5N$ is mandatory for models with small $c$.
Models with large $c$ saturate earlier and are less sensitive to this
constraint.

The effective $\beta$-function is estimated by central differences,
\begin{equation}
  \beta_{\rm eff}(N_i) = \frac{\Omega_{i+1} - \Omega_{i-1}}{\ln N_{i+1} - \ln N_{i-1}},
\end{equation}
and plotted against $\Omega_{\rm late}(N_i)$ at the central point.
All simulations use exact statevector evolution ($N \leq 12$, dense trace)
or typicality averaging ($N \geq 14$, $10$ random states).

% ─────────────────────────────────────────────────────────────────────────────
\section{Results}
\label{sec:results}

\subsection{Autonomous Logistic Flow}

Figure~\ref{fig:collapse} shows $\beta_{\rm eff}/c$ vs.\ $\Omega_{\rm late}$
for eight models with monotonic flow.
The data collapse onto the universal curve $-\Omega(1-\Omega)$ with
slope $= 1.003$, intercept $= -0.004$, $\chi^2/\text{dof} = 1.11$,
and $p = 0.31$, confirming that the logistic form
$\beta(\Omega) = -c\,\Omega(1-\Omega)$
is the leading-order description of the flow.

The collapse is non-trivial: data from $N = 4$--$20$ across models
spanning two orders of magnitude in $c$ lie on a single curve.
This demonstrates that $\beta_{\rm eff}$ is \emph{autonomous}---the rate
of change of information accessibility depends only on the current state
of accessibility, independently of the microscopic scale at which the
system is probed.
To our knowledge, this is the first empirical characterization of an
autonomous accessibility flow extracted from quantum hardware.

Three lines of evidence independently support the logistic form over the
power-law alternative $\beta = -\alpha\Omega$:
(i) $\Delta\text{AIC} = -4.72$ favoring logistic for the KI chaotic
9-point series;
(ii) the curvature parameter $\gamma(N) = d\ln\Omega/d\ln N$ steepens
from $-3.8$ to $-6.8$ with increasing $N$ for SYK ($q = 4, 6$),
consistent with the logistic UV fixed point becoming active;
(iii) the log-log plot $\ln\Omega$ vs.\ $\ln N$ shows convex curvature
for SYK ($\Delta R^2 = 0.022$), which the logistic produces
and the power-law does not.

\begin{figure}[h]
  \centering
  \includegraphics[width=\columnwidth]{figures/fig1_collapse.pdf}
  \caption{Normalized effective $\beta$-function $\beta_{\rm eff}/c$ vs.\
    $\Omega_{\rm late}$ for eight models. Data from $N = 4$--$20$ collapse
    onto the universal curve $-\Omega(1-\Omega)$ (solid line), confirming
    the autonomous logistic flow. Slope $= 1.003$,
    $\chi^2/\text{dof} = 1.11$, $p = 0.31$.}
  \label{fig:collapse}
\end{figure}

\subsection{Continuous IR Fixed-Point Spectrum}

Table~\ref{tab:spectrum} lists $c \pm \sigma(c)$ for all eleven models.
The spectrum $c \in [0.17, 4.30]$ is continuous and structured:

\emph{Non-chaotic regime} ($c < 0.4$, non-monotonic $\Omega$):
KI $J = 0.3, 0.5$. The flow is non-monotonic; fits are unreliable.
These models are excluded from the collapse analysis.

\emph{Luttinger class} ($c \approx 1$--$1.3$):
KI $J = 0.9$ ($c = 1.260 \pm 0.050$) and $J = 1.0$ ($c = 0.993 \pm 0.020$).
The value $c \approx 1$ is consistent with the $c_{\rm CFT} = 1$ central
charge of the Luttinger liquid universality class.

\emph{Generic chaotic} ($c \approx 1.9$--$3.1$):
KI $J = 0.6$--$0.8$. The spectrum is non-monotonic in $J$, with a maximum
at $J \approx 0.8$ ($c = 3.116 \pm 0.111$). KI $J = 0.8$ nearest-neighbor
approaches SYK all-to-all, demonstrating that $c$ does not measure
interaction range.

\emph{Schwarzian/JT class} ($c \approx 3.9$--$4.3$):
SYK $q = 4$ ($c = 3.901$) and $q = 6$ ($c = 4.300$).
The near-independence of $c$ on $q$ is consistent with SYK reaching the
Schwarzian/JT universality class in the IR, though $c$ is a finite-size
scaling exponent distinct from the Schwarzian coupling~\cite{Kitaev2015,Maldacena2016}.

Figure~\ref{fig:spectrum} displays the complete spectrum.

\begin{table}[h]
\caption{Scaling exponent $c$ for all eleven models.
  $\sigma(c)$ from nonlinear least-squares covariance matrix.
  Non-monotonic models excluded from collapse analysis.}
\label{tab:spectrum}
\begin{ruledtabular}
\begin{tabular}{lcccc}
Model & $c$ & $\sigma(c)$ & $R^2$ & IR class \\
\hline
KI $J=0.3$ & 0.169 & — & 0.623 & Non-chaotic$^\dagger$ \\
KI $J=0.5$ & 0.356 & — & 0.787 & Non-chaotic$^\dagger$ \\
KI $J=1.0$ & 0.993 & 0.020 & 0.994 & Luttinger \\
KI $J=0.9$ & 1.260 & 0.050 & 0.996 & Luttinger \\
KI $J=0.6$ & 1.916 & — & 0.933 & Chaotic \\
KI $J=0.7$ & 2.719 & 0.039 & 0.991 & Chaotic \\
KI $J=0.8$ & 3.116 & 0.111 & 0.970 & Chaotic (max) \\
KI $1/r$   & 3.406 & 0.477 & 0.982 & Predicted$^\ddagger$ \\
SYK $q=4$  & 3.901 & 0.238 & 0.9996 & Schwarzian/JT \\
KI mixing  & 4.170 & — & — & Inverse flow \\
SYK $q=6$  & 4.300 & 0.191 & 0.9997 & Schwarzian/JT \\
\end{tabular}
\end{ruledtabular}
\footnotetext[1]{$\dagger$ Non-monotonic $\Omega$; fit unreliable.}
\footnotetext[2]{$\ddagger$ Pre-registered prediction confirmed;
  see Sec.~\ref{sec:prediction}.}
\end{table}

\begin{figure}[h]
  \centering
  \includegraphics[width=\columnwidth]{figures/fig2_spectrum.pdf}
  \caption{Continuous IR fixed-point spectrum $c$ for eleven models,
    ordered by increasing $c$. Error bars show $\pm\sigma(c)$ from
    covariance matrix. Dashed lines indicate the Luttinger class
    ($c \approx 1$) and Schwarzian/JT class ($c \approx 3.9$).}
  \label{fig:spectrum}
\end{figure}

\subsection{Pre-Registered Prediction Confirmed}
\label{sec:prediction}

Prior to executing the simulation, we predicted that a model with
intermediate-range interactions---between nearest-neighbor KI and
all-to-all SYK---would yield $c \in [1, 4]$.
For KI with power-law coupling $J(r) = J_0/r$ (all pairs $i < j$,
$J_0 = 1.0$, $N = 4$--$12$), we obtain $c = 3.406 \pm 0.477$, $R^2 = 0.982$.

The result $c = 3.41$ falls within the predicted range, between
$c_{\rm KI\,NN} \approx 1.0$ (nearest-neighbor) and $c_{\rm SYK} \approx 3.9$
(all-to-all), and is physically coherent: long-range interactions produce
faster scrambling than nearest-neighbor but do not reach maximal.
This model was not used in constructing the hypothesis; its confirmation
constitutes an independent test of the framework.

Figure~\ref{fig:prediction} shows $\Omega_{\rm late}(N)$ for the
power-law model alongside the fit.

\begin{figure}[h]
  \centering
  \includegraphics[width=\columnwidth]{figures/fig3_prediction.pdf}
  \caption{$\Omega_{\rm late}(N)$ for KI power-law coupling $J(r) = 1/r$
    (circles), with algebraic fit $c = 3.406 \pm 0.477$ (solid line).
    The result falls within the pre-registered prediction $c \in [1, 4]$
    (shaded band), between nearest-neighbor KI ($c \approx 1$, dashed)
    and SYK ($c \approx 3.9$, dotted).}
  \label{fig:prediction}
\end{figure}

\subsection{Observable Independence}

To test whether $c$ is a property of the Hamiltonian or of the measurement
protocol, we repeated the KI chaotic $J = 1.0$ simulation with two
choices of scrambled operator: $V = Z_1$ ($r = 1$) and $V = Z_2$ ($r = 2$).

The results are:
$c(Z_1) = 0.9939$, $c(Z_2) = 1.0062$, difference $1.2\%$.
$A(Z_1) = 0.6417$, $A(Z_2) = 0.5068$, difference $21\%$.

The exponent $c$ is invariant under change of observable, while the
amplitude $A$ absorbs the dependence on operator distance.
This separation between universal exponent and non-universal amplitude
is the standard signature of renormalization group universality: the
critical exponent is independent of the ``regularization'' (choice of
observable), while the amplitude carries the non-universal information.

% ─────────────────────────────────────────────────────────────────────────────
\section{Discussion}
\label{sec:discussion}

\subsection{Non-Monotonic $c(J)$ and the Luttinger Point}

The spectrum $c(J)$ is non-monotonic: it rises from $c \approx 0$ at
$J = 0.3$ to a maximum $c \approx 3.1$ at $J = 0.8$, then falls to
$c \approx 1$ at $J = 1.0$.
The value $J = 1.0$ is a special point in the KI phase diagram, consistent
with the Luttinger liquid universality class ($c_{\rm CFT} = 1$).
Two models ($J = 0.9$ and $J = 1.0$) share this class, suggesting a
Luttinger basin of attraction in the IR.

\subsection{UV Symmetry Class Independence}

We tested whether $c$ depends on the Altland-Zirnbauer (AZ) symmetry
class of the SYK model.
For $k = 2N$ Majorana modes, the SYK model alternates between AZ classes
as $N$ increases ($k \bmod 8 = 0$ for $N = 4, 8, 12$; $k \bmod 8 = 4$
for $N = 6, 10, 14$).
Fitting separately: $c_A = 3.97 \pm 0.19$, $c_B = 3.81 \pm 1.99$,
difference $4\%$.
The large error in $c_B$ (two effective fit points) prevents a definitive
conclusion, but the result is consistent with UV symmetry class independence.

Combined with observable independence, three tests confirm that $c$ is
an IR property insensitive to UV structure:
observable choice ($1.2\%$), UV symmetry class ($4\%$), and J-value
(varies continuously---only changing the full Hamiltonian changes $c$).

\subsection{Log-Log Curvature as Evidence for Logistic Flow}

In a log-log plot, a pure power-law $\Omega \propto N^{-\alpha}$ appears
as a straight line.
The logistic form $\Omega = 1/(1 + AN^c)$ introduces convex curvature
as $N$ grows, because the UV fixed point ($\Omega \to 1$) contributes
a $(1-\Omega)$ factor absent in the power-law.
For SYK $q = 4$: $R^2_{\rm linear} = 0.976$, $R^2_{\rm quadratic} = 0.998$
($\Delta R^2 = 0.022$); for SYK $q = 6$: $\Delta R^2 = 0.022$.
This curvature is a third independent line of evidence for the logistic
form, complementing $\Delta\text{AIC}$ and $\gamma(N)$ steepening.
Figure~\ref{fig:scaling} illustrates both effects.

\begin{figure}[h]
  \centering
  \includegraphics[width=\columnwidth]{figures/fig4_scaling.pdf}
  \caption{Main: $\Omega_{\rm late}(N)$ for KI chaotic $J = 1.0$
    ($N = 4$--$20$, circles, log-log scale). Solid: logistic fit
    ($c = 0.993$, $A = 0.642$, $\Delta\text{AIC} = -4.72$). Dashed:
    power-law. Arrow marks $N \approx 45$ where fits diverge by
    $>15\%$. Inset: $\ln\Omega$ vs.\ $\ln N$ for SYK $q = 4$;
    linear fit (dashed) vs.\ quadratic fit (solid),
    $\Delta R^2 = 0.022$, demonstrating convex curvature consistent
    with the logistic UV fixed point.}
  \label{fig:scaling}
\end{figure}

\subsection{Limitations}

\emph{N$\leftrightarrow\mu$ identification:}
The identification of system size $N$ as the RG scale $\mu$ is an
emergent analogy supported by~\cite{Kim2025,Vidal2007,Swingle2012}
but not derived from first principles.
$\beta_{\rm eff}$ is defined operationally; its interpretation as a
Wilsonian $\beta$-function is a working hypothesis.

\emph{Hip.~A vs.\ Hip.~C discrimination:}
The logistic and power-law forms are statistically indistinguishable
for $\Omega \ll 0.1$ (deep IR, large $N$).
The $\Delta\text{AIC} = -4.72$ for KI chaotic depends primarily on the
data range $\Omega > 0.1$.
Definitive discrimination requires $N \approx 45$ (logistic predicts
$\Omega \approx 0.031$; power-law predicts $\Omega \approx 0.037$,
a $>15\%$ difference), accessible via MPS/DMRG or cluster HPC.

\emph{SYK $c$ values:}
The values in Table~\ref{tab:spectrum} for SYK are obtained from
disorder-averaged fits; direct fits to averaged data give
$c \approx 4.55$ (q=4) and $c \approx 5.37$ (q=6).
The difference arises because the nonlinear fit to averaged data $\neq$
average of per-seed fits. Table values from per-seed averaging are
used as authoritative throughout.

% ─────────────────────────────────────────────────────────────────────────────
\section{Conclusion}
\label{sec:conclusion}

We have demonstrated that the late-time OTOC saturation value
$\Omega_{\rm late}(N)$ defines an autonomous renormalization-group-like
flow $\beta(\Omega) = -c\,\Omega(1-\Omega)$ that is universal in form
across dynamical classes and model-dependent in its single coefficient $c$.

The exponent $c$ constitutes a continuous, observable-independent
empirical spectrometer of the effective IR fixed point, resolving
eleven models into a structured spectrum with physically interpretable
classes: Luttinger liquid ($c \approx 1$), generic chaotic
($c \approx 2$--$3$), and Schwarzian/JT ($c \approx 4$).
A pre-registered prediction was confirmed with an independent model,
and $c$ was shown to be invariant under changes of local observable and
UV symmetry class---the standard signature of RG universality.

These results establish $c$ as an experimentally accessible, hardware-independent
probe of dynamical universality class in quantum many-body systems.
The relationship between $c$ and the surface gravity $\kappa$ of effective
horizons, suggested by the observed $c \propto \ln\kappa$ scaling in
thermal variants of these models, is explored in a companion paper.

Full simulation code, raw data, and analysis scripts are available
at~\cite{PerezEugenio2026c}.

\begin{acknowledgments}
The author used AI-assisted tools for manuscript preparation, simulation
execution, and logical consistency checks.
The views expressed are those of the author.
\end{acknowledgments}

% ─────────────────────────────────────────────────────────────────────────────

\bibliography{kaelion_v8}

\end{document}
